Modern AI systems combine capabilities such as tool use, structured output, streaming, reasoning, and multimodal processing. Failures may emerge only from particular capability combinations, creating stochastic interaction bugs that are difficult to isolate through independent testing. Exhaustive combinatorial discovery can identify such interactions but becomes increasingly expensive as capability count and interaction order grow.

We present SCIF v0.0.4, a residual-risk-guided method for stochastic mixed-order interaction discovery. SCIF first screens and confirms pairwise interactions, suppresses those already supported by the data, measures the remaining empirical failure risk, and invokes higher-order localization only when substantial unexplained risk remains.

Evaluation on the synthetic BSIB benchmark family produced 6,999 exact recoveries across 7,000 unseen holdout runs (99.9857\%), with no observed false interaction candidates. In a 100-seed method comparison, SCIF V4 recovered all evaluated pairwise, overlapping, pure-three-way, and mixed-order structures; pairwise SCIF and deletion localization both achieved 0\% exact recovery on the mixed scenario. For that eight-capability benchmark, SCIF required 22,800 mean simulator executions versus 93,000 for exhaustive order-three discovery, a 75.48\% reduction.

Ablation experiments confirmed distinct roles for residual-risk detection and higher-order localization. Across scaling experiments from eight to twenty capabilities, SCIF achieved 20/20 observed exact recoveries at each size while execution reduction reached 96.14\%. These results support residual-risk-guided escalation as a promising approach for stochastic interaction discovery, while external validation beyond synthetic benchmarks remains future work.
